\section{Height Gap and Curve Identification}\label{sec:interface}

The field results identify harmonic conditional means away from prescribed
strands. To identify the curve, we must determine their boundary traces on
the strand and retain those traces when the strand is stopped. This section
separates the corresponding implications. Every application below requires
the uniform versions of the preceding field and geometric estimates; a
statement for one fixed conditioning would not suffice.

\subsection{Conditional means viewed from harmonic landings}

Throughout this section $\gamma_n$ denotes the ordered \emph{dual sign
strand}, rather than the affine zero segments. Write $\mathcal F^n_t$
for the sigma-field generated by its oriented prefix up to time $t$.
It reveals adjacent signs, rather than the height ratios encoded by
affine crossing positions. Let $T$ be a stopping time for this filtration.
For a deterministic smooth test $f$ define
\[
 X_n(f)=\int_{D_n}f(z)h_n(z)\dd z,\qquad
 M^n_T(f)=\E[X_n(f)\mid\mathcal F^n_T].
\]
Thus $M^n_t(f)$ is an exact martingale. A test chosen afresh after $T$ does
not define this martingale and will not be used in the driving argument.
For a complete strand, write
$m_n(z)=\E[h_n(z)\mid\gamma_n]$.
The same notation with a superscript $T$ denotes the stopped conditional
mean. Comparisons with bounded pinned data bound these profiles uniformly.
Their subsequential limits are harmonic on the complementary open sets:
Proposition~\ref{b:cell} supplies the interior weak equation,
Proposition~\ref{b:response} gives $A>0$, and the common-wall comparison
supplies compactness.

Let $v_n$ be an interior observer vertex, with a fixed positive physical
distance from the original boundary. An independent walk from $v_n$
samples the harmonic measure
of the complete strand and the original boundary. A landing on the
positive side is understood as an oriented access; two prime ends with the
same Euclidean location need not have the same sign. We condition on that
side and, when necessary, on the finitely many possible orientations of
the incident lattice edge. All statements about boundary samples below are
made under this harmonic sampling law. Lemma~\ref{g:extra} is needed to
exclude additional excursions near the incident strand.

For the global boundary-sampling statements, an experiment means a
Dobrushin domain with sign-compatible exterior heights
$+\lambda$ and $-\lambda$, where $\lambda$ is fixed as $n\to\infty$.
For any finite collection of experiments fix
$0<\lambda_0\le\min\lambda\le1/2$. Constants in estimates whose
barriers attach to original boundary may depend on $\lambda_0$.
The artificial bounded-data conditionings used for local comparison
are not additional global experiments to which boundary avoidance
is asserted.

Here is the exclusion required before using the interior records of
(G7)--(G8). Let $\xi_n$ be the walk landing and let
$d_n=\operatorname{dist}(v_n,\partial D_n)$ in lattice units.
Lemmas~\ref{g:boundary-landing} and~\ref{g:observer-landing} give
\begin{equation}\label{h:landing-exclusions}
 \lim_{\eta\downarrow0}\limsup_{n\to\infty}
 \Prob\left(\xi_n\hbox{ lies on a strand side},\
   \operatorname{dist}(\xi_n,\partial D_n\cup\{v_n\})<\eta d_n
       \right)=0.
\end{equation}
The boundary estimate is the consequence of exterior-attaching barriers
and the independent full-walk exit comparison proved there; it does
not use SLE or a matched height gap. Its stronger full-walk form
also applies to landings on an arbitrary stopped prefix.

When conditioning on a side or on an incident-edge label, first take
a subsequence along which its probability converges. If the limit is
zero, its entire contribution to a bounded harmonic test vanishes.
Otherwise its probability is bounded below along that subsequence,
so dividing \eqref{h:landing-exclusions} by that lower bound gives
the same exclusion for the conditioned sampling law. There are only
finitely many such labels at each stage. This is the only
normalization needed to pass the exclusion to the boundary samples
used below.

\begin{lemma}[Harmonic sampling in the physical slit]\label{h:slit-sampling}
Let the original Jordan lattice domains be contained in one fixed
bounded set. Let $U_n$ be the component containing an observer $z_n$
after removing a simple dual prefix, or after cutting along a
complete simple dual crosscut. Suppose
$\operatorname{dist}(z_n,\partial U_n)\ge\rho>0$ in physical units.
Label its prime-end boundary by the original positive and negative
arcs and the two oriented strand sides; empty arcs are allowed.
There are at most four cyclic arcs before completion and at most
two in a completed component.

Kill triangular-lattice walk on its first visit to the original
boundary or to a vertex adjacent to the explored strand. Original
boundary values take priority at a vertex of both types.
Let $\xi_n$ denote this landing with its incident-side label.
There is a coupling with Brownian exit $\xi_n^{\rm B}$ from $U_n$
such that, for every $a>0$,
\[
 \sup_{U_n,z_n}
 \Prob\{|\xi_n-\xi_n^{\rm B}|>a
           \text{ or their labels differ}\}\longrightarrow0 .
 \tag{H2}
\]
The supremum includes all such prefixes, hence all stopping rules.
Repeated Euclidean boundary points are distinguished by their
incident accesses.
\end{lemma}
\begin{proof}
We first recall the required deterministic screening argument.
For a simply connected polygonal slit $U$, two complementary
prime-end arcs, and $R>0$, let $L_R$ be the union of internal
paths of diameter at most $R$ joining those arcs. The frontier
of the observer component of $U\setminus L_R$ is contained in
two balls of radius $2R$. Indeed, classify a crosscut by which
of the two arc endpoints it separates from the observer.
Two crosscuts in one class touching the observer component
must intersect, and their diameter bound then puts that entire
class in one ball. The argument uses prime ends on a slit,
not an identification of its two sides. This is the planar
argument in \cite[equation (3.55)]{SSdiscrete}.
Beurling's estimate therefore gives
\[
 \Prob_z\{\text{Brownian motion enters }L_R\}
       \le C(R/\rho)^\alpha .
\]
Apply this to each arc against its complementary arc. For a
partition into $m$ cyclic arcs, entrance into the union of
short paths joining different labels has probability at most
$Cm(R/\rho)^\alpha$, for a fixed $\alpha>0$.

The lattice walk has the required physical incidence. Every
primal edge crossed by an explored dual edge has both endpoints
observed and absorbing. Thus the interpolated walk cannot cross
the dual strand before it is killed. An interior killing vertex
has a connector of diameter $C/n$ through its own dual hexagon
to an incident strand edge, on the side bearing its sign.
If several strand edges are incident, every such access has
that same sign. At original boundary the connector is trivial
and the incoming edge specifies its access. Consecutive explored
dual edges cross primal edges belonging to a common triangle;
their absorbing endpoints form a connected set attached to the
original absorbing boundary. This set lies within $C/n$ of
the physical boundary and reaches across the annuli used below.
It supplies the connected obstacle for the discrete Beurling
estimate, including at the last explored dual vertex.

Fix $n^{-1}\ll r\ll R\ll\rho$. Couple the unrestricted
interpolated walk, with its diffusive clock, and Brownian motion
so that their distance up to a fixed time cutoff is less than
$r/4$, except with probability $o_n(1)$. The ordinary planar
invariance principle gives this coupling independently of the
slit; the triangular walk is isotropic. Both exit-time tails
at the cutoff are bounded by those for a fixed enclosing ball.
For each process use its \emph{own} first entrance into the
$2r$ boundary layer. Beurling's estimate at these two marginal
stopping times confines each remaining tail to diameter $R$,
except with total probability $C((r+n^{-1})/R)^\alpha$.
No Markov assertion about the joint coupling filtration is made.

Let $\theta$ be the first time either coupled path is within
$r$ of the physical boundary. On successful coupling both
paths are still alive, both of their $2r$ layer times have
occurred, and their positions at $\theta$ can be joined by a
segment inside $U_n$. If the exit labels differ, their confined
tails, this segment, and the lattice terminal connector give
an internal path of diameter $CR$ joining different labels.
The Brownian path has therefore entered the corresponding
mixed-access set. The preceding screening bound gives
\[
 \Prob\{\text{label mismatch or exit distance}>CR\}
 \le o_n(1)+C(r/R)^\alpha+Cm(R/\rho)^\alpha
       +\text{cutoff error}.                              \tag{H3}
\]
Take the time cutoff large, then $R$ small, then $r/R$ small,
and finally $n$ large. This proves (H2).
In particular the proof compares the walk directly with the
physical dual slit, not with an unmarked Hausdorff approximation
or a thickened domain with possibly many label changes.
\end{proof}

\begin{lemma}[Marked limits and local access]\label{h:marked-limits}
In the setting of Lemma~\ref{h:slit-sampling}, let the observers
converge, $z_n\to z$, and take a Hausdorff subsequence
of the closed removed sets consisting of the explored slit and the
original exterior. Retain the conformal maps from the disk to the
observer components, normalized by $f_n(0)=z_n$ and $f_n'(0)>0$,
and the inverse images of all boundary-label transition accesses.
After a further subsequence on which these marks converge,
bounded harmonic functions with the prescribed finite
label values converge on compact subsets to the Poisson
extension of the limiting marked arc data. Their harmonic
boundary samples converge with these marks. At a surviving
common local boundary piece of two such domains, the limit
preserves the incident-side identification.
\end{lemma}
\begin{proof}
The protected observer and the common enclosing ball make the
normalized inverse maps a nondegenerate normal family. Kernel
convergence identifies their limit with the observer component.
Extract also all transition angles on the unit circle, preserving
their cyclic order. The associated bounded step functions converge
at every angle except the finitely many limiting transitions.
Poisson integration proves compact convergence; normalized
arclength gives zero mass to those transitions. Colliding
transitions merely delete an arc of vanishing harmonic mass.
This retains the boundary marks rather than reconstructing
them from the limiting closed set.

For approach samples, use a free Brownian path independent of
all domain and field records. At fixed positive distance from
the boundary, kernel convergence and compact harmonic convergence
give convergence of its smoothed readings. The boundary-layer
tail estimate in the preceding proof lets that distance decrease
to zero after taking the mesh limit. Equivalently, the normalized
conformal maps give the limiting Brownian exit angle.

Finally suppose the two finite domains and their labels agree
on the component of a ball of radius $R$ next to a common
oriented strand. Their bounded label-Poisson extensions differ,
at a point in the concentric ball of radius $R/2$ and at distance
$r\ll R$ from that strand, by at most
$CM(r/R)^\alpha$: the difference is zero on the common local
boundary, and Beurling bounds escape from the ball.
The same bound holds if the common boundary is an original
arc. It is uniform in the domains, so it persists in marked
kernel limits. The mixed-access screening in the preceding
proof excludes conflating distinct sides in this procedure.
This proves the asserted local identification, without
claiming continuity of prime-end labels from closed sets alone.
\end{proof}

For the global experiments, the original domain approximations supply
the common bounded enclosure. Before applying these lemmas to
complete-strand samples, fix $\rho>0$ below the observer's original
boundary clearance and discard the bounded contribution from
configurations whose observer-to-strand distance is less than $\rho$
in physical units. Lemma~\ref{g:observer-landing} makes this contribution
vanish as $\rho\downarrow0$ after the mesh limit. On the retained
configurations both lemmas apply with this fixed clearance. For a
positive-probability side or orientation branch, divide the discarded
probability by its lower bound as above. This localization estimates
reading errors under the original sampling law; the lower comparison
in (G7) remains on all its admissible records.

For a bounded harmonic limiting profile $m$ and an independent Brownian
motion $W$ in its component, $m(W_t)$ has an almost sure limit as $W$
approaches its exit. Denote that limit by $Y$. It is a random boundary
value, not yet a deterministic height gap. The following elementary lemma
specifies the comparison needed to make it deterministic.

\begin{lemma}[A common lower bound for two-scale readings]
\label{h:common-kernel}
For $1\le i\le k$, let $U^i_n,V^i_n$ take values in $[-M,M]$, with
$U^i_n$ measurable with respect to a sigma-field $\mathcal A^i_n$.
Suppose that
\[
 \E|U^i_n-V^i_n|\longrightarrow0.
\]
Suppose there are probability measures $\nu_n$ on $[-M,M]$, a constant
$c>0$, and events $G^i_n\in\mathcal A^i_n$ of probability tending to one,
such that, on $G^i_n$,
\[
 \Prob(V^i_n\in B\mid\mathcal A^i_n)\ge c\,\nu_n(B)
 \quad\text{for every Borel set }B.
\]
Then every subsequence has a further subsequence on which all the
$U^i_n,V^i_n$ converge in probability to the same deterministic constant.
\end{lemma}
\begin{proof}
Pass to a subsequence on which $\nu_n$ and the laws of all $U^i_n$ converge.
Conditional integration gives
\[
 \E|U^i_n-V^i_n|
 \ge c\,\E\left[\mathbf1_{G^i_n}
          \int|U^i_n-y|\,\nu_n(\dd y)\right].
\]
The lost contribution outside $G^i_n$ is at most $2M\Prob((G^i_n)^c)$.
If $\mu_i$ and $\nu$ are the limiting laws, bounded continuity yields
$\int|x-y|\,\mu_i(\dd x)\nu(\dd y)=0$.
A product of two probability measures supported on the diagonal must
consist of the same point mass in each coordinate. The point is common
to all $i$ because $\nu$ is common. Agreement of the readings completes
the proof.
\end{proof}

The constant $c$ in this lemma need not approach one. It must, however,
stay positive as unused scale buffers grow. A different reference law for
each external configuration would not provide the stated conclusion.

\begin{lemma}[Locality of bounded conditional means]\label{h:locality}
Assume the sequentially uniform comparison and trace estimates of
Theorem~\ref{c:comparison} and Theorem~\ref{t:trace}, and the local
weak equation of Proposition~\ref{b:cell}. Consider two conditioning
systems with uniformly bounded mean profiles and the same lattice
domain and original fixed data inside a ball of radius $R$. Suppose
they share a connected oriented strand across that ball and all sign
observations inside it. For averaged observations at distance of order $r$
from that strand, supported in the concentric ball of radius $R/2$,
where $r\to\infty$ in lattice units and $r/R\to0$,
the two conditional means differ by a quantity tending to zero.

There is also an exterior version. If the two systems agree outside
a ball of radius $\rho$, their original exterior data agree, and their
unchanged strands connect that ball to radius $R$, then at distance
of order $r$ from its center, with $\rho\ll r\ll R$, the limiting
mean difference is at most $CM(\rho/r)^{1/2}$. Here $M$ bounds the
individual profiles. In both assertions use translates and dilates
of the fixed radial probability kernel from
Section~\ref{sec:trace-analysis}. A sampling ball centered at $y$
has radius $s$ less than one tenth of both $r$ and the distance from
$y$ to the boundary of the common free region. For each fixed $a>0$,
retain $s\ge ar$ while taking the mesh limit in coordinates rescaled
by $r$; the cutoff may then decrease to zero.
\end{lemma}
\begin{proof}
For the first assertion compare the two fields with common reflecting
sign walls. On shared strands their limiting mean difference has zero
Brownian-tested trace by the comparison and trace estimates. The
common original fixed data give zero differences at the corresponding
pins, so the same trace argument gives zero trace on the common
original boundary pieces. Away from these observed boundaries the
difference is harmonic by the weak equation and $A>0$. Its
absolute value on the other boundary pieces is at most $2M$.
The maximum principle therefore bounds it by $2M$ times the
probability of escaping the common ball before hitting the connected
common strand. Beurling bounds that probability by $C(r/R)^{1/2}$.
For arbitrary sequences with $r\to\infty$ and $r/R\to0$, rescale
by $r$ and first restrict to a common ball of radius $Lr$ about a
nearest strand point, with $L$ fixed. Centrality puts this ball inside
the shared region for all sufficiently large scales. Local $L^1$
compactness transfers its $CM L^{-1/2}$ bound to the fixed-cutoff
averages. First take the mesh limit, then $L\to\infty$, and finally
remove the averaging cutoff. Sequential uniformity gives the
assertion for every admissible conditioning sequence.

For the second assertion apply the identical argument in the common
exterior region, with the small changed window as the only uncontrolled
boundary. The probability of reaching that window from distance $r$
before hitting a connected common arm is at most $C(\rho/r)^{1/2}$.
The bounds persist under averaging in a sufficiently small ball about
the observation point. In both assertions the zero trace is a
conclusion of the common-wall comparison, not an equality of the
sampled wall heights.
\end{proof}

\begin{proposition}[The two readings required for the height gap]
\label{h:readings}
Assume the harmonic-mean compactness and common-sign trace comparison of
Proposition~\ref{b:cell}, Theorem~\ref{c:comparison}, and
Theorem~\ref{t:trace}. Assume also the landing exclusions
\eqref{h:landing-exclusions} and the sampled-configuration estimates
of Proposition~\ref{g:sample-comparison} and Lemma~\ref{g:extra}.
Then, along any joint subsequence of finitely many experiments, the
following construction gives the hypotheses of
Lemma~\ref{h:common-kernel}.
The limiting variables are the Brownian boundary samples of the
conditional mean profiles.
\end{proposition}
\begin{proof}
Fix the positive-probability side and orientation branches described
above. We construct an inner reading $V_n$ from the incident strand
and an outer reading $U_n$ from the external arms. Both approximate
the same boundary sample. The comparison of sampled configurations
will then bound the conditional law of $V_n$ below by one common
reference law, independently of the external record.

For a target error, first use \eqref{h:landing-exclusions} to discard
strand landings within distance $\eta d_n$ of the original boundary
or observer, choosing $\eta>0$ so that their probability is small.
Recenter a retained landing at zero. Choose the lattice radii
\[
 1\ll r_{\rm in}\ll R_{\rm core}\ll R_{\rm in}\ll R_{\rm mid}
       \ll R_{\rm conn}\ll r_{\rm out}\ll n.                 \tag{H1}
\]
Require also $r_{\rm out}<\eta d_n/20$. Thus the largest disks used by
(G7)--(G8), including their factors $6$ and $4$, lie inside the original
domain and exclude the observer. Clearance is measurable in the outer
record, which includes the domain and observer; retaining it imposes no
additional condition on the inner strand. Write
$\mathcal A_n=\sigma(\Phi_{R_{\rm mid}})$ under the specified landing
event $Z^\sigma$, and let $G_n\in\mathcal A_n$ be the retained
clearance event. The diagonal choices below make $\Prob(G_n)\to1$.

\emph{Agreement of complete-profile readings.}
First take mesh limits at fixed positive physical approach distances.
In each limiting component the bounded harmonic martingale has a
boundary limit. Its values at first entrances into balls of radii
tending to zero about the eventual landing have that same limit:
these entrance times approach the exit time pathwise. They need not
be stopping times after the future-dependent center is specified.

At an interior approach point use a radial smooth average in a ball
smaller than one tenth of its distance to the strand and of the
approach radius, with the fixed probability kernel defining $P_s$ in
Section~\ref{sec:trace-analysis}. The limiting harmonic average equals
the value at its center. Impose a positive lower cutoff on the averaging radius
before taking the mesh limit, then remove it. At each fixed positive
physical approach scale the sampling point is interior almost surely,
so this cutoff loses arbitrarily little probability. The comparison
compactness and the vanishing original--filled $L^1$ error give
convergence of the conditional means in $L^1$ on every compact subset
of the limiting free regions, on a joint subsequential realization.
At a fixed radius cutoff $a>0$, the sampling supports have distance
at least $9a$ from the strand; the retained landing clearance and (H1)
also keep them a fixed positive distance from the original exterior.
Hausdorff convergence therefore places all these supports in one
compact subset $C_a$ of the limiting free regions. Since the kernels
have supremum norm at most $\|\varphi\|_\infty a^{-2}$, the discrepancy
of their readings is bounded uniformly by
$\|\varphi\|_\infty a^{-2}\|m_n-m\|_{L^1(C_a)}$, which tends to zero.
This transfers both readings to the lattice, including for kernels
selected from the records.
A diagonal choice of the radii in (H1) therefore makes the readings of
the complete profile at $r_{\rm in}$ and $r_{\rm out}$ agree in $L^1$.
The scales are chosen slowly; no convergence rate is asserted.

\emph{The reading determined by the inner record.}
Retain the incident oriented strand and terminal walk piece in
$\Theta_{R_{\rm in}}$. Define its canonical mean in the disk of radius
$2R_{\rm in}$, with zero exterior heights and precisely those oriented
strand observations. This conditioning is feasible by assigning
sufficiently small signed heights. The construction is the same
measurable function of the inner record in every experiment.

Apply (G8) at $(R_{\rm core},R_{\rm in})$ and average its conditional
bound. Except on an event of vanishing probability under $Z^\sigma$,
the walk does not leave $R_{\rm in}$ after its first entrance at
$R_{\rm core}$. Its first entrance at $r_{\rm in}$ therefore belongs
to the recorded terminal piece, and reversing that piece determines
the sampling point. Any complete-chord piece in $B(0,R_{\rm core})$
other than the incident strand $\beta_{R_{\rm in}}^\sigma$ would
give two additional crossings of the same annulus. Outside the same
type of exceptional event, the canonical and complete conditionings
thus agree throughout $B(0,R_{\rm core})$.

Choose the inner radial kernel using the distance to the recorded
strand and the cutoff just described. Its canonical mean reading,
denoted $V_n$, is determined by $\Theta_{R_{\rm in}}$ alone.
If this record does not contain the required approach point, or its
kernel cutoff fails, set $V_n=0$. This is the same bounded function
of the inner record in every experiment; it does not use $G_n$.
Lemma~\ref{h:locality}, with
$r_{\rm in}/R_{\rm core}\to0$, makes its difference from the
complete-profile inner reading tend to zero. The extra-excursion
estimate supplies both the common strand neighborhood and the
presence of the chosen approach point in the recorded walk.

\emph{The reading determined by the outer record.}
Retain the external arms and initial walk in $\Phi_{R_{\rm mid}}$.
Delete every oriented dual-edge observation whose adjacent primal
sign pair meets $B(0,R_{\rm conn})$; retain the remaining external-arm
edge observations and both vertices of each pair. Every retained sign
vertex still has an oppositely constrained neighbor.

Apply (G8) at $(R_{\rm mid},R_{\rm conn})$ and average its bound.
Except on an event of vanishing probability under $Z^\sigma$, the
unrecorded connection of the two external arms stays inside
$B(0,R_{\rm conn})$: the arms already supply two disjoint annular
crossings, and an excursion of their connection would supply two more.
Hence the modified and complete conditionings agree outside
$B(0,R_{\rm conn}+C)$, with an absolute lattice constant $C$.
All retained observations are determined by $\Phi_{R_{\rm mid}}$.
The first walk entrance at $r_{\rm out}$ is also in that record,
since it precedes the first entrance at $R_{\rm mid}$.

Use a radial kernel smaller than one tenth of $r_{\rm out}$ and
of the distance to the recorded external strands, with a positive
lower cutoff as above. Its support stays outside the changed window.
The mean under the modified conditioning, tested with this kernel,
is $\mathcal A_n$-measurable; denote it by $U_n$. Set $U_n=0$
when its required outer clearance or kernel cutoff fails.
In the common exterior component, limiting mean differences are bounded
harmonic functions with zero trace on the unchanged strands and original
boundary. The only uncontrolled boundary is the deleted window, where
the difference has absolute value at most $2M$. A connected common arm
crosses the intervening annulus. The maximum principle and Beurling's
estimate therefore bound the influence of this window at the sampling
ball by
\[
 C M(R_{\rm conn}/r_{\rm out})^{1/2}.
\]
The bounded enlargement by $C$ is negligible at these scales.
Sequentially uniform compactness and trace comparison transfer the
bound to the lattice with an error tending to zero. Combining the two
localizations with complete-profile agreement gives
$\E|U_n-V_n|\to0$.

\emph{One common lower kernel and the order of limits.}
For each $n$, fix one admissible reference external record at radius
$4R_{\rm in}$. Equation (G7) bounds every inner-record law below by
$C^{-1}$ times this reference law, for every admissible external record
at that radius in every experiment. Since
$R_{\rm mid}>4R_{\rm in}$, the record at $R_{\rm mid}$ is a
restriction of the record at $4R_{\rm in}$. Averaging the bound over
the latter conditional on $\Phi_{R_{\rm mid}},Z^\sigma$ preserves
$C^{-1}$. Push this one reference law through the canonical inner-mean
function defining $V_n$ to obtain $\nu_n$. Lattice symmetries identify
the finite list of landing orientations. The exceptional events used
to make the readings agree do not restrict this lower comparison, and
no factor is paid for unused annuli.

Here is the scale order for a finite collection of limiting Brownian
boundary readings and a requested error less than $\eps>0$.
Choose the clearance exclusion before any physical buffers. Choose the
outer physical approach scale for both the boundary-reading error and
the retained clearance. Next choose $R_{\rm conn}/r_{\rm out}$
for exterior locality and $R_{\rm mid}/R_{\rm conn}$ for its
excursion error. Choose $R_{\rm in}\ll R_{\rm mid}$, then
$R_{\rm core}/R_{\rm in}$ for the inner excursion error, and finally
$r_{\rm in}/R_{\rm core}$ for inner locality and the inner boundary
reading. These choices are in physical units before the mesh limit, so
every retained lattice scale diverges. Hold the scales and kernel
cutoffs fixed while taking that limit, and then remove the cutoffs.
Repeat for $\eps\downarrow0$ and take a diagonal subsequence.
The comparison constant in (G7) is independent of all these unused buffers.
Lemma~\ref{h:common-kernel} now applies.
\end{proof}

\begin{proposition}[Subsequential universal height gap]\label{h:height-gap}
Assume the uniform inputs of Proposition~\ref{h:readings} and the
trace bounds of Lemmas~\ref{g:sign-trace} and~\ref{g:upper-half}.
Every mesh subsequence has a
further subsequence and a deterministic $l\in(0,1/2]$ such that all
terminal conditional mean profiles in any prescribed finite collection
of the sign-compatible Dobrushin experiments specified above have
Brownian boundary trace $+l$ on a
positive strand side and $-l$ on a negative strand side.
\end{proposition}
\begin{proof}
Include a fixed reference Dobrushin experiment with original boundary
heights $\pm1/2$. The barrier estimate gives nonzero harmonic mass to
its path sides. Compactness supplies joint mean profiles and harmonic
samples. Proposition~\ref{h:readings} and
Lemma~\ref{h:common-kernel} make the sampled positive boundary value a
single deterministic constant, shared by all included experiments.
The reference experiment's positive lower trace bound gives $l>0$,
and its complete-crosscut comparison gives $l\le1/2$.
Sign reflection identifies the negative value as $-l$.

Apply this argument to a countable collection of interior balls, with
a diagonal subsequence. If a terminal profile had a set of incorrect
boundary values of positive harmonic measure from any ball in that
collection, its sampled boundary value would have positive probability
of differing from $l$. Thus the statement holds almost surely in
every component seen from that exhaustion. Original boundary arcs retain
their prescribed values by the boundary comparison, rather than by the
height-gap argument.
\end{proof}

\subsection{Passing to a stopped strand}

For an independent walk from $z$, let $\mu^z_{n,T}$ be harmonic measure on
the boundary of the component of $D_n\setminus\gamma_n[0,T]$ containing
$z$. The nonreturn input needed here is the following consequence of
Proposition~\ref{g:nonreturn}: for each compact set of protected observers,
\begin{equation}\label{h:nr}
 \lim_{s\downarrow0}\limsup_{n\to\infty}\sup_T
 \E\int_{\gamma_n[0,T]}
 \mathbf1_{\{\operatorname{dist}(x,\gamma_n(T,\mathrm{end}])<s\}}
       \,\mu^z_{n,T}(\dd x)=0.
\end{equation}
This is an annealed harmonic-mass estimate, uniformly in the stopping
rule. It does not assert that every realized stopped prefix has a
small conditional return probability.

\begin{proposition}[Stopped harmonic observable]\label{h:stopped}
Fix a subsequence furnished by Proposition~\ref{h:height-gap} and its
deterministic height gap $l$. In this proposition $n$ tends to infinity
only along that subsequence. Assume \eqref{h:nr}, the survival bound
of Lemma~\ref{g:survival}, and the sequentially uniform field
comparisons. Harmonic sampling and boundary marks are as in
Lemmas~\ref{h:slit-sampling} and~\ref{h:marked-limits}.
Let $H^l_{n,T}$ be the continuum harmonic
function in the stopped physical domain with the original exterior
values and values $+l,-l$ on the new sides. For every deterministic
smooth test $f$ supported in a protected interior ball,
\begin{equation}\label{h:stopped-observable}
 \sup_T\E\left|M^n_T(f)-\int f(z)H^l_{n,T}(z)\dd z\right|
 \longrightarrow0.
\end{equation}
The supremum is over stopping rules before a fixed neighborhood of
that ball is reached, and the assertion holds in either direction
of exploration.
\end{proposition}
\begin{proof}
It suffices to treat an arbitrary sequence $T_n$, since failure of the
supremum would select a counterexample sequence. Use one independent
full walk $S$. Write $\xi$ for its old-prefix landing, $\upsilon$
for its completed-strand landing, and set
\[
 \mathcal G_n=\sigma(\gamma_n[0,T_n],S),\qquad
 A_n=\{\xi=\upsilon\},\qquad
 q_n=\Prob(A_n\mid\mathcal G_n).
\]
For any $\mathcal G_n$-measurable $Z_n\in[0,2M]$,
\begin{equation}\label{h:reweight}
 \E Z_n\le 2M\Prob(q_n<p)+p^{-1}\E[Z_n\mathbf1_{A_n}].
\end{equation}
Indeed, on $\{q_n\ge p\}$ integrate the identity
$\E[Z_n\mathbf1_{A_n}\mid\mathcal G_n]=Z_nq_n$.
Lemma~\ref{g:survival} makes the first term arbitrarily small
with a fixed $p>0$, uniformly in $T_n$. This is the required
change of sampling law; mere absolute continuity is not used.

At physical radius $r$ let $V^{\rm old}_{n,r}$ be a small radial
average of the stopped conditional mean at the first walk entrance
into $B(\xi,r)$. Its averaging radius is less than one tenth of
the distance to the observed strand and one tenth of $r$.
Define $V^{\rm end}_{n,r}$ in the same way with the complete strand
and $\upsilon$. Positive cutoffs on the averaging radius are removed
after the mesh limit, as in Proposition~\ref{h:readings}.
Let $b(\xi)$ be $+l$ or $-l$ according to the old oriented side,
and put
\[
 Z_n=|V^{\rm old}_{n,r}-b(\xi)|
                \mathbf1_{\{\xi\notin\partial D_n\}} .
\]
This is measurable with respect to $\mathcal G_n$.
By Proposition~\ref{h:height-gap} and
Lemmas~\ref{h:slit-sampling}--\ref{h:marked-limits},
\[
 \lim_{r\downarrow0}\limsup_n
 \E\!\left[|V^{\rm end}_{n,r}-b(\upsilon)|
                \mathbf1_{\{\upsilon\notin\partial D_n\}}\right]=0.
 \tag{H4}
\]
This is the bounded harmonic martingale boundary limit in the
completed component, sampled by the lattice walk through (H2).
First discard arbitrarily small completed-observer clearances
using Lemma~\ref{g:observer-landing}; (H2) then applies with a
fixed protected observer.
The first entrance around its eventual landing need not itself
be a stopping time: these approach points tend to the exit
pathwise, which is all that this boundary-limit assertion uses.

Discard old-side landings within distance $2R$ of the original
boundary and those within distance $2R$ of the future strand.
Their total probability, denoted $e_n(R)$, satisfies
$\lim_{R\downarrow0}\limsup_n e_n(R)=0$ by
\eqref{h:landing-exclusions} and \eqref{h:nr}.
On the remaining event, and on $A_n$, the two readings use the
same walk segment, the same label, and the same averaging kernel
when $r<R/10$. The prefix extends across this common ball:
its starting point is on the distant original boundary and its
tip belongs to the distant future, or to the original boundary
if exploration is complete. The stopped and completed
conditionings have exactly the same sign observations there.
Lemma~\ref{h:locality} therefore bounds their mean difference by
$CM(r/R)^\alpha+o_n(1)$ at fixed $r,R$. Consequently
\[
 \E[Z_n\mathbf1_{A_n}]
 \le \E\!\left[|V^{\rm end}_{n,r}-b(\upsilon)|
              \mathbf1_{\{\upsilon\notin\partial D_n\}}\right]
       +CM(r/R)^\alpha+2M e_n(R)+o_n(1).                 \tag{H5}
\]
First choose the survival error in \eqref{h:reweight}, obtaining
$p>0$. Choose $R$ so that the discarded mass is small compared
with $p$ times the desired final error. Then choose $r\ll R$
to make both (H4) and the locality error equally small.
Keep these physical scales and the kernel cutoffs fixed while
taking $n\to\infty$. This order proves that the old-side boundary
reading equals its prescribed label in $L^1$, uniformly in
the chosen sequence of stopping rules.

Original boundary accesses retain their original values.
For clarity, the argument proving (GB2)--(GB3) also yields
\[
 \lim_{R\downarrow0}\limsup_n
 \Prob\{d_{D_n}(S_{\tau_{D_n}},\gamma_n)<R\}=0
 \tag{H6}
\]
in physical units: condition on the independent original-boundary
exit, apply pure-arc repulsion away from the opposite arc, and
use the two-ball bound at the two sign-change accesses.
On the complementary event the final approach has a neighborhood,
in its original-domain component, free of the completed strand.
The original-boundary comparison gives the prescribed trace there.
The pathwise coupling used for (H2), including confinement of the
terminal tails, transfers this internal-access assertion to
Brownian motion. Thus (H6) applies also on the event that the
old landing lies on the original boundary. No boundary avoidance
deduced from SLE is used.

Finally extract jointly the stopped closed sets, bounded conditional
profiles, normalized conformal maps, and their boundary marks.
The limiting stopped profile is harmonic by the weak equation.
The preceding readings give its boundary values at Brownian-almost
every marked access. Bounded harmonic uniqueness and
Lemma~\ref{h:marked-limits} identify this profile with the limit
of $H^l_{n,T_n}$. The comparison compactness and original--filled
$L^1$ estimate give convergence of the lattice conditional means
against the fixed smooth test $f$. The harmonic functions $H^l_{n,T_n}$
converge uniformly on its protected support. Integration against $f$
and boundedness therefore prove \eqref{h:stopped-observable}.
Proposition~\ref{f:averages} supplies the moments of the underlying
field averages if those averages are retained in the joint limit.

Reverse exploration reveals the same kind of sign prefixes and
satisfies the same comparisons and nonreturn estimate. Repeating
the proof gives the reverse assertion.
\end{proof}

\begin{remark}[No transfer from absolute continuity alone]
A sequence of equivalent measures can give vanishing mass to a set
whose mass stays positive under another access law. For example,
on $[0,1]$ let $\nu$ be Lebesgue measure and give $\mu_n$ density $2/n$
on $[0,1/2]$ and $2(1-1/n)$ on its complement. Then
$\mu_n([0,1/2])\to0$ while $\nu([0,1/2])=1/2$.
The survival reweighting and nonreturn estimate supply the uniform
local transfer in the preceding proof. The exact tower identity
for conditional means does not replace either estimate.
\end{remark}

\subsection{Driving convergence and the uniform curve metric}

First obtain the deterministic $l$ on a subsequence using the reference
experiment with exterior heights $\pm1/2$. Now add the experiment with
exterior heights $\pm l$ and apply Proposition~\ref{h:height-gap} to
these two experiments on that subsequence. Any further extracted
height gap must still equal $l$, since the reference readings
already converge to $l$. This yields the stopped observable for
the matched experiment without assuming its existence before
$l$ was fixed.

On this height-gap subsequence we use original boundary heights
$\pm l$, matching the value on the
new sides. Map the marked domain to the upper half-plane, with the
starting point at zero and the target at infinity. If $g_t$ is its
capacity-parametrized Loewner map and $U_t$ its driver, the harmonic
observable is, with a fixed choice of side orientation,
\begin{equation}\label{h:angle}
 H^l_t(z)=l-\frac{2l}{\pi}\arg(g_t(z)-U_t).
\end{equation}
There is no boundary force-point term because the old and new heights
match.

\begin{lemma}[The observable input for driving convergence]
\label{h:driving}
For simple lattice crosscuts in marked Jordan domains, suppose that
the exact conditional-mean martingales of deterministic protected
interior averages satisfy \eqref{h:stopped-observable}, uniformly over
the stopping rules of the driving argument, with matched exterior
heights $\pm l$, where $l>0$. Suppose the comparison is to the continuum
slit-domain function \eqref{h:angle}. Then the forward chordal drivers
converge, locally uniformly in capacity time in law, to $2B$.
\end{lemma}
\begin{proof}
We specify which part of the harmonic-observable criterion is used.
Let $x_t<y_t$ be the two images of the starting point under $g_t$.
For $t>0$ introduce the bounded-driver coordinates
\[
 s=\log(y_t-x_t),\qquad
 G_s(z)=\frac{2g_t(z)-x_t-y_t}{y_t-x_t},\qquad
 \widetilde U_s=\frac{2U_t-x_t-y_t}{y_t-x_t}\in[-1,1].
\]
Let $\mathcal F_s$ be the sign-prefix filtration at the corresponding
time, equivalently the filtration of the explored dual curve after
this deterministic coordinate change. In these coordinates, as in
\cite[Sections 4.2--4.5]{SSdiscrete},
fix a finite interval $[-S,S]$ and the two deterministic interior
probes protected throughout that interval. Replace each point probe by
a radial smooth average in a sufficiently small protected disk.
The harmonic mean-value property makes the average of the continuum
observable equal to its value at the center. Thus
\eqref{h:stopped-observable} supplies the same two approximate
martingales as the point probes.

For stopping times $s_0\le s_1$ with
$s_1-s_0\le\eps^2$ and driver oscillation at most $\eps$, the two-probe
Taylor calculation gives, outside an event of arbitrarily small
prescribed probability,
\begin{align}
 \left|\E[\Delta\widetilde U+
       2\widetilde U_{s_0}\Delta s\mid\mathcal F_{s_0}]\right|
       &\le C_S\eps^3, \label{h:driver-first}\\
 \left|\E[(\Delta\widetilde U)^2-
       2(1-\widetilde U_{s_0}^{\,2})\Delta s
           \mid\mathcal F_{s_0}]\right|
       &\le C_S\eps^3. \label{h:driver-second}
\end{align}
Here $\Delta s=s_1-s_0$ and
$\Delta\widetilde U=\widetilde U_{s_1}-\widetilde U_{s_0}$.
To obtain these estimates from qualitative convergence, fix
$S,\eps$ and the exceptional probability first, and only then take
the mesh sufficiently fine. Conditional Markov inequality converts
the $L^1$ observable errors into the needed conditional errors.
No predetermined rate is required.

To apply the approximate-diffusion criterion, start the comparison
diffusion at the same bounded-driver value as the discrete process.
For a fixed comparison horizon $S$, choose the exceptional tolerance
$\eta<\delta^5/S^2$, where $\delta$ is the criterion's approximation
parameter, before choosing the mesh; Proposition~4.5 of
\cite{SSdiscrete} then gives its comparison on $[0,S-1]$.
For a limiting window $[-S,S]$, start instead at $-S'$ and let
$S'\to\infty$ after the mesh limit. The coupling argument in the
proof of Theorem~4.1 of \cite{SSdiscrete}, pp.~115--116, removes the
initial value and identifies the stationary whole-line
bounded-coordinate diffusion. Thus an arbitrary initial Jacobi law
is not substituted for the chordal limit.

The deterministic Taylor argument and the approximate-diffusion
argument of \cite[Propositions 4.2 and 4.5]{SSdiscrete} use precisely
these martingale and harmonic-observable inputs; no Gaussian
conditional law enters these steps. They identify the bounded-coordinate
diffusion with squared diffusion coefficient $2(1-u^2)$ and drift
$-2u$. The inverse coordinate change \cite[Section~4.6]{SSdiscrete} gives
the chordal driver $2B$.
The discrete harmonic-measure comparison in that source is replaced
here by the explicit continuum comparison assumed in
\eqref{h:stopped-observable}.
\end{proof}

For completeness, the normalization of this identification can also be
seen directly in any continuous Loewner subsequential limit. If the
angle at one protected point is a continuous martingale, the identity
\[
 U_t=\operatorname{Re}g_t(z)-
       \operatorname{Im}g_t(z)\cot\arg(g_t(z)-U_t)
\]
makes the driver a semimartingale after localization. Writing
$U=V+N$, where $V$ has finite variation and $N$ is a continuous local
martingale, and $Z_t=g_t(z)-U_t$, It\^o's formula gives the finite-variation
part of its angle as
\[
 -\operatorname{Im}(Z_t^{-1})\dd V_t+
 \operatorname{Im}(Z_t^{-2})
       \left(2\dd t-\tfrac12\dd[N]_t\right).
\]
Two protected points with different
$\operatorname{Re}(Z_t^{-1})$ force $\dd V_t=0$ and
$\dd[N]_t=4\dd t$. A dense set of probes provides such a pair locally.
L\'evy's characterization then gives $U=2B$.
This calculation identifies a continuous limit; the
approximate-diffusion argument is still needed to obtain driving
tightness.

\begin{proposition}[Uniform convergence of the whole curve]
\label{h:uniform-path}
If the hypotheses of Lemma~\ref{h:driving} hold for both exploration
directions, the interfaces converge to chordal $\mathrm{SLE}_4$ in
uniform Euclidean distance modulo increasing reparametrization.
\end{proposition}
\begin{proof}
The dual interfaces are simple crosscuts and are continuously driven
in each direction. Forward and reverse driving convergence follows
from Lemma~\ref{h:driving}. Corollary 1.8 of \cite{SSun} applies to these
two separate convergence-in-law statements for $\kappa\le4$; it does
not require a previously constructed joint limit of the two drivers.
Its conclusion is convergence of the entire curve, including its
endpoints, in the uniform curve metric on the compactified reference
domain.

Uniform convergence of the marked Jordan boundary parametrizations
to the smooth limit boundary gives uniform convergence of the
normalized inverse conformal maps on the closed reference disk
by Rad\'o's theorem \cite[Section 1, p.~182]{Rado}.
One way to see the boundary part is to use uniform local connectivity
of these parametrized Jordan boundaries to turn small reference
crosscuts into uniformly small physical crosscuts; interior convergence
is the kernel theorem. Thus uniform curve convergence transfers to
Euclidean distance in $D$. If the meaning of boundary convergence
were weakened to kernel convergence alone, this last inference would
need an additional hypothesis.

In every traversed triangular face, the affine zero segment and the
corresponding ordered dual strand stay within $C/n$ of each other.
Their order agrees, so an increasing reparametrization gives the same
$C/n$ bound for the complete curves. Hence the choice of polygonal
representative does not alter the conclusion.
\end{proof}

\subsection{The variance normalization}

The height-gap argument so far allows $l$ to depend on the extracted
subsequence. The conditional covariance identifies it.

\begin{proposition}[Identification of the gap]\label{h:normalization}
Assume the joint conclusions of Proposition~\ref{h:stopped},
Theorem~\ref{t:covariance}, and the field limit
Theorem~\ref{b:gff}, with the uniform moments needed to pass conditional
second moments. Then every subsequential height gap satisfies
\[
 l^2=\frac{\pi\sigma^2}{8}.
\]
\end{proposition}
\begin{proof}
Use the experiment with original heights $\pm l$, whose interface limit
is $\mathrm{SLE}_4$ by the preceding propositions. Stop before compact
smooth tests approach the curve. The conditional means are the
corresponding harmonic observables, while their conditional covariance
is $\sigma^2$ times the Dirichlet Green kernel of the remaining domain.
The sequential uniformity of Theorem~\ref{t:covariance} is used to
apply it to the actual random stopped strings. The moments from
Proposition~\ref{f:averages} justify convergence of the averaged
conditional variances.

In the half-plane set $Z_t(z)=g_t(z)-2B_t$. It\^o's formula yields
\[
 \dd\log Z_t(z)=-\frac{2}{Z_t(z)}\dd B_t.
\]
The cross variation of the harmonic observables is therefore
\begin{equation}\label{h:mean-bracket}
 \frac{16l^2}{\pi^2}
       \operatorname{Im}\frac1{Z_t(z)}
       \operatorname{Im}\frac1{Z_t(w)}\dd t.
\end{equation}
For covariance $(-\Delta)^{-1}$ the half-plane Green function is
\[
 G_{\mathbb H}(z,w)=\frac1{2\pi}
       \log\left|\frac{z-\overline w}{z-w}\right|.
\]
Differentiating its conformal pullback gives
\begin{equation}\label{h:green-loss}
 -\partial_tG_{D_t}(z,w)=\frac2\pi
       \operatorname{Im}\frac1{Z_t(z)}
       \operatorname{Im}\frac1{Z_t(w)}.
\end{equation}
These normalizations and the observable-bracket/Green-loss comparison
agree with \cite[preprint version, Lemmas 2.5--2.6,
equations (2.3)--(2.4)]{SScontinuum}.
Only this calculation is used; the stronger conditional coupling
in that work is not an input.

Integrate both identities against a nonzero nonnegative smooth test
$f$ twice. No pointwise value of the limiting field is used. The law
of total variance equates the gained variance of the conditional-mean
martingale with the expected loss of the residual Green variance.
For a nonzero localized time interval this gives
\[
 \frac{16l^2}{\pi^2}\,
 \E\int\left(\int f(z)\operatorname{Im}(1/Z_t(z))\dd z\right)^2\dd t
 =
 \frac{2\sigma^2}{\pi}\,
 \E\int\left(\int f(z)\operatorname{Im}(1/Z_t(z))\dd z\right)^2\dd t.
\]
The common factor is finite and strictly positive. Cancelling it
proves the formula.
\end{proof}

\begin{corollary}[Universality]\label{h:assembly}
For the model of Section~\ref{sec:model}, the constant
\[
 \lambda_*=\sigma\sqrt{\pi/8}
\]
belongs to $(0,1/2]$, is independent of the marked domain, and gives
the curve conclusion of Theorem~\ref{thm:main}. The field stiffness
is $vA=\sigma^{-2}$, and its limiting mean is the harmonic extension
of the prescribed exterior values.
\end{corollary}
\begin{proof}
Apply the reference height-gap extraction to every subsequence.
Proposition~\ref{h:normalization} gives the same value $\lambda_*$
on each extraction; its membership in $(0,1/2]$ comes from the
reference complete-crosscut bounds. Include, in each finite comparison,
the fixed experiment with exterior heights $\pm\lambda_*$.
Every subsequence then has a further subsequence to which the driving
and uniform-path conclusions apply with the same matched value.
This proves full convergence. The field assertions are the
identifications in Theorem~\ref{b:gff} and
Proposition~\ref{b:response}.
\end{proof}
